% EM structure and teaching revision 20260928
\subsection{为何有环路定律仍要判断对称性}
静磁场满足 $\oint\boldsymbol B\cdot\dd\boldsymbol l=\mu_0I_{\mathrm{enc}}$。只有在所选路径上磁场方向和大小受到足够对称性约束时，才能从一个积分方程直接求出局部场。
无限长圆柱导体半径为 $R$、电流密度均匀，则圆形安培回路给出
\[
 B_\phi(r)=\begin{cases}\mu_0Ir/(2\pi R^2),&r<R,\\
 \mu_0I/(2\pi r),&r>R.
 \end{cases}
\]
导体内部包围的是部分电流，故磁场随 $r$ 增大；外部才是反比律。若同一圆柱中的电流密度不均匀且不具有轴对称性，环路定律仍成立，但不能把 $B$ 直接提到积分号外。
磁场的散度为零则约束另一个积分：任意闭合面的净磁通量为零。它不要求任意开曲面的磁通量为零，也不禁止磁力线穿过该曲面。计算线圈磁通时必须指定开曲面、面积正向和匝数；线圈的磁链是各匝磁通之和，不一定等于一匝磁通乘统一匝数。

